由 (5)，恒等式 (6) 化为
\[
\begin{array}{@{}l@{\quad}l@{}}
(7)& \displaystyle G(\varphi(x);F_1(\varphi_1\cdots\varphi_n);\cdots;F_n(\varphi_1\cdots\varphi_n)) =H(\varphi_1(x)\cdots\varphi_n(x))=0;
\end{array}
\]
由于 $\varphi_1(x)\cdots\varphi_n(x)$ 在代数上无关，(7) 不仅关于
$x$ 恒等成立，而且关于各 $\varphi_i$ 也恒等成立；也就是说，关于
独立参数 $\lambda$ 有恒等式：
\[
\begin{array}{@{}l@{\quad}l@{}}
(8)& \displaystyle H(\lambda_1\cdots\lambda_n)=G(\mu_1\lambda_1+\cdots+\mu_n\lambda_n;F_1(\lambda_1\cdots\lambda_n);\cdots;F_n(\lambda_1\cdots\lambda_n))=0.
\end{array}
\]
由恒等式 (8) 可以推出，方程
\[
\begin{array}{@{}l@{\quad}l@{}}
(9)& \displaystyle f(x)=x^n-F_1(\lambda_1\cdots\lambda_n)x^{n-1}+F_2(\lambda_1\cdots\lambda_n)x^{n-2}\cdots\pm F_n(\lambda_1\cdots\lambda_n)=0
\end{array}
\]
\emph{相对于} $\Omega(\lambda_1\cdots\lambda_n)$ \emph{具有群}
$\Gamma$。事实上，这个恒等式表明，可以有理地给出的
$\lambda=\mu_1\lambda_1+\cdots+\mu_n\lambda_n$ 是一个属于
$\Gamma$ 的根函数；而由代换 $\lambda_i=\varphi_i(x)$ 可见，不可能
再有一个属于 $\Gamma$ 的某个子群的函数也能有理地给出；同时，在同一
代换下，$\lambda$ 还与它的所有共轭量互不相同。此外，若 $\Omega^*$
是任意一个包含 $\Omega$ 的数域，则 $f(x)$ \emph{相对于}
$\Omega^*(\lambda_1\cdots\lambda_n)$ \emph{具有群} $\Gamma$；
因为由第 2 节可知，在作代换 $\lambda_i=\varphi_i(x)$ 时，即使添入
$\Omega^*$，群也不会约化。

现在还须证明，在第 1 节所修正的意义下，$f(x)$ 确实遍历所有
$g_\Gamma$。为此要注意，由 (5) 可知，$F_i(\lambda)$ 是关于其自变量
$\lambda$ 的代数无关函数。因此，方程组
\[
\begin{array}{@{}l@{\quad}l@{}}
(10)& \displaystyle F_1(\lambda_1\cdots\lambda_n)=-a_1;\quad F_2(\lambda_1\cdots\lambda_n)=a_2;\quad \cdots;\quad F_n(\lambda_1\cdots\lambda_n)=\pm a_n
\end{array}
\]
对于任意值组 $a_1\cdots a_n$ 都有解，只有下文尚待指出的奇异值组
例外。对于相对于 $\Omega^*$ 的某个 $g_\Gamma$ 所对应的每个值组
$a_1\cdots a_n$，相应的 $\lambda_i$ 作为属于该群的自然无理量，
\srcfn{*)}{另见下文！} 都是 $\Omega^*$ 中的量。因此，当
$\lambda_1\cdots\lambda_n$ 遍历所有值组时，$g(x)$ 就遍历所有方程
（在修正后的意义下）；其中包含对应于 $\Omega^*$ 中各值的全部
$g_\Gamma$。反过来，根据 Hilbert 不可约性定理，$\Omega^*$ 中的
$\lambda$ 值\glqq{}一般地\grqq{}也对应于方程 $g_\Gamma$；所以，
在第 1 节所修正的意义下，参数表示 (9) 确实给出了相对于 $\Omega^*$
的全部 $g_\Gamma$。

现在只剩下确定使参数表示失效的奇异值组。将最小基中的函数拆成分子
和分母，写为
\[
\varphi_i(x)=\frac{\psi_i(x)}{\chi_i(x)}.
\]
把它们代入恒等式 (5)，设这些恒等式（不作约分）化为：
\[
\begin{array}{@{}l@{\quad}l@{}}
(11)& \displaystyle \sigma_k(x)=\frac{G_k(x)}{H_k(x)}\quad\text{或}\quad H_k(x)\cdot\sigma_k(x)=G_k(x),
\end{array}
\]
于是所有 $H_k(x)$ 的乘积可以被所有分母 $\chi_i(x)$ 的乘积整除。
对于所有满足
\[
H_1(\alpha)\cdot H_2(\alpha)\cdots H_n(\alpha)\neq0
\]
的根组 $\alpha_1\cdots\alpha_n$，所有分母 $\chi_i(\alpha)$ 也都
不消失，因而可以在 (11) 中除以 $H_k(\alpha)$，得到
\[
\begin{array}{@{}l@{\quad}l@{}}
(12)& \displaystyle \sigma_k(\alpha)=\frac{G_k(\alpha)}{H_k(\alpha)}=F_k(\varphi_1(\alpha);\cdots;\varphi_n(\alpha)),
\end{array}
\]
其中，由于分母不消失，各 $\varphi_i(\alpha)$ 都取有限值；对于方程
$g_\Gamma$，这些值还是 $\Omega^*$ 中的量。只要确定
$\alpha_1\cdots\alpha_n$，使 $a_k=(-1)^k\sigma_k(\alpha)$，(12)
就给出 (10) 的解系。\srcfn{*)}{由此，方程的求解便归结为求解由独立
函数组成的方程组：
\[
\varphi_1(\alpha_1\cdots\alpha_n)=\lambda_1;\quad \cdots;\quad \varphi_n(\alpha_1\cdots\alpha_n)=\lambda_n.
\]}

因此，只有满足
\[
H_1(\alpha)\cdot H_2(\alpha)\cdots H_n(\alpha)=0;
\]
的系数组 $a_k=(-1)^k\sigma_k(\alpha)$ 才可能成为\emph{奇异的}；
对这些系数组，(11) 不再给出上述表示，而变成 $0=0$。这时必须另作
研究，才能判定这样定义的奇异 $a_1\cdots a_n$ 中是否包含对应于方程
$g_\Gamma$ 的值组，以及它们对应于哪些数域。\srcfn{**)}{例如，正如
F. Seidelmann 在前引处所证明的，这种对应于奇异值组的方程只在三次
和四次方程的交错群情形下出现，而且只对那些包含三次单位根域的数域
$\Omega^*$ 出现。每种情形都可以给出一个简单的补充表示。} 综上可得

\emph{定理：属于群 $\Gamma$ 的 Lagrange 型域若有一个最小基}
\srcfn{***)}{从一个最小基的存在立刻可以推出任意多个最小基的存在；
它们都能由原来的最小基通过可逆有理变换导出。}
\emph{，其系数属于 $\Omega$，那么对于每个包含 $\Omega$ 的数域
$\Omega^*$，相对于 $\Omega^*$、具有预先指定的群 $\Gamma$ 的全部
方程，都可以通过参数表示 (2) 作有理构造——可能例外的是相对于该参数
表示为奇异的方程，但这些方程可以另行给出。参数必须遍历 $\Omega^*$
中所有不导致群约化的值。若最小基的系数域 $\Omega$ 等于所有有理数
所成的域 $\mathfrak R$，则对于任意数域都可以作这种构造。}

根据 Hilbert 不可约性定理，从参数流形中排除 $\Omega^*$ 内那些导致
群约化的值，并不会降低该流形的维数。因此还有：
\emph{最小基的存在推出，相对于 $\Omega^*$、具有群 $\Gamma$ 的方程
所成流形的维数，等于无特约方程所成流形的维数；特约并不降低这一
维数。}

